\begin{abstract}
Gravitational coupling between a protoplanetary disc and an embedded eccentric planet is an important, long-standing problem, which has been not yet been conclusively explored. Here we study the torque and associated orbital evolution of an eccentric planet in a two-dimensional disc via the semi-analytical, fully global linear approach. Our methodology has the advantage that the spatial structure of the density waves launched by the planet is solved for fully. This allows us to account for the possibility of torque excitation over an extended radial interval for each Fourier harmonic of the perturbation, as opposed to earlier approximate treatments localized around Lindblad and corotation resonances. We systematically explore the torque behaviour across the space of disc properties (assuming power law profiles for the disc surface density and temperature), including the aspect ratio. Crucially, we examine the torque variation as the orbital eccentricity becomes supersonic relative to the gas motion (when planetary eccentricity is of order the disc aspect ratio), finding that the torque robustly reverses its sign near this transition. We also find that for shallow surface density gradients planetary migration may become outwards beyond this transition, although the rapid eccentricity damping (which is typically $\sim 10^2$ times faster than the orbital migration rate) would quickly restore inwards migration as the planet circularizes. Our self-consistently computed torques are in qualitative agreement with past numerical studies of eccentric planet-disc coupling. We provide our torque data for different disc parameters to the community for future testing and implementation in population synthesis studies.
\end{abstract}